\section*{Примечания о чтениях}
\subsection*{\texttt{GELFAND-P001415}}
\par\noindent\leavevmode\phantomsection\label{GELFAND-0154-U001}\textbf{\texttt{GELFAND-0154-U001}} (eq16-cancelled).\nopagebreak[4]\par
Под (16) зачёркнута строка «а при r … 0». Перечёркивающая линия пересекает знак отношения.\par
Кандидаты: а при r≠0; а при r≥0. Чтение не установлено в дипломатическом исходнике.\par
Историческое примечание исходной транскрипции: Оригинал 2408×3448; координаты (x0,y0,x1,y1)=(335, 880, 740, 950), начало в левом верхнем углу. Просмотр увеличенного фрагмента: 5×. Зачёркивание сохранено; знак не восстанавливался из условия рекурсии.\par\medskip
\par\noindent\leavevmode\phantomsection\label{GELFAND-0154-U002}\textbf{\texttt{GELFAND-0154-U002}} (continued-fraction-cancelled-word).\nopagebreak[4]\par
После «W\_\{-r\}(τ) опр.» находится короткий фрагмент, закрытый несколькими густыми косыми штрихами.\par
Кандидаты: не записаны. Чтение не установлено в дипломатическом исходнике.\par
Историческое примечание исходной транскрипции: Оригинал 2408×3448; координаты (x0,y0,x1,y1)=(375, 1230, 1060, 1390), начало в левом верхнем углу. Просмотр увеличенного фрагмента: 3×. Сохраняется отдельный маркер зачёркнутого фрагмента. Надёжное буквенное чтение невозможно.\par\medskip
\par\noindent\leavevmode\phantomsection\label{GELFAND-0154-U003}\textbf{\texttt{GELFAND-0154-U003}} (eq19-right-hand-side).\nopagebreak[4]\par
В (19) отчётливо написано W\_r(τ)W\_\{1-r\}(τ)=1.\par
Кандидаты: W\_r(τ)W\_\{1-r\}(τ)=1. Чтение не установлено в дипломатическом исходнике.\par
Историческое примечание исходной транскрипции: Оригинал 2408×3448; координаты (x0,y0,x1,y1)=(155, 1600, 1810, 2315), начало в левом верхнем углу. Просмотр увеличенного фрагмента: 2×. Возможная ошибка источника: при формальном подставлении выражений (20)–(21) произведение равно q/(q+1)\textasciicircum{}2, а не 1 (для общих допустимых q и индексов). Это только проверка согласованности, не поправка к тексту. Единица сохранена.\par\medskip
\par\noindent\leavevmode\phantomsection\label{GELFAND-0154-U004}\textbf{\texttt{GELFAND-0154-U004}} (w-index).\nopagebreak[4]\par
Повторяющийся индекс W и показатель q имеют короткую начальную черту и загнутый вправо нижний штрих; явной точки над ними нет.\par
Кандидаты: r; i. Чтение не установлено в дипломатическом исходнике.\par
Историческое примечание исходной транскрипции: Оригинал 2408×3448; координаты (x0,y0,x1,y1)=(375, 1230, 1060, 1390), начало в левом верхнем углу. Просмотр увеличенного фрагмента: 3×. В наборе предварительно использовано r; индексы i в (23) имеют иной, точечный рисунок. Не заменять эту палеографическую неопределённость математическим предположением.\par\medskip
\par\noindent\leavevmode\phantomsection\label{GELFAND-0154-U005}\textbf{\texttt{GELFAND-0154-U005}} (eq23-overwritten-exponents).\nopagebreak[4]\par
В правом верхнем столбце числителя (23) малые индексы после i покрыты повторными штрихами; далее читается +n−1.\par
Кандидаты: i\_n+n−1; i\_4+n−1. Чтение не установлено в дипломатическом исходнике.\par
Историческое примечание исходной транскрипции: Оригинал 2408×3448; координаты (x0,y0,x1,y1)=(165, 2330, 2200, 3448), начало в левом верхнем углу. Просмотр увеличенного фрагмента: 1.5×. Последняя строка даёт более ясное i\_n+n−1. Это служит внутренним сравнением, но не снимает неопределённости двух исправленных верхних индексов.\par\medskip
\subsection*{\texttt{GELFAND-P001416}}
\par\noindent\leavevmode\phantomsection\label{GELFAND-0154-U006}\textbf{\texttt{GELFAND-0154-U006}} (mayer-name).\nopagebreak[4]\par
В скобках под заголовком читается «Линдон Майер (?)»; вопросительный знак находится в самой рукописи.\par
Кандидаты: Линдон Майер (?); Линден Майер (?). Чтение не установлено в дипломатическом исходнике.\par
Историческое примечание исходной транскрипции: Оригинал 2408×3448; координаты (x0,y0,x1,y1)=(230, 960, 715, 1050), начало в левом верхнем углу. Просмотр увеличенного фрагмента: 3×. Первое слово не отождествлялось с внешними библиографическими сведениями. На следующем кадре видна самостоятельная запись «D. Mayer»; обе формы сохраняются.\par\medskip
\par\noindent\leavevmode\phantomsection\label{GELFAND-0154-U007}\textbf{\texttt{GELFAND-0154-U007}} (interval-graph).\nopagebreak[4]\par
Рисунок в прямоугольной рамке, приблизительно квадратный. Видны четыре восходящих наклонных отрезка, несколько близких к вертикали штрихов, повторные обводки и короткие отдельные штрихи; подписей нет.\par
Кандидаты: не записаны. Чтение не установлено в дипломатическом исходнике.\par
Историческое примечание исходной транскрипции: Оригинал 2408×3448; координаты (x0,y0,x1,y1)=(1410, 1330, 1710, 1630), начало в левом верхнем углу. Просмотр увеличенного фрагмента: 4×. Нужна проверка точной связности в верхних и нижних точках и статуса коротких штрихов. Не предполагаются ни стрелки, ни дополнительные значения осей. Схема не перерисована.\par\medskip
\par\noindent\leavevmode\phantomsection\label{GELFAND-0154-U008}\textbf{\texttt{GELFAND-0154-U008}} (slope-letter).\nopagebreak[4]\par
В знаменателе после первой суммы и перед верхним индексом −1 написана петлеобразная буква с длинным нижним штрихом и индексом j.\par
Кандидаты: ρ\_j; p\_j. Чтение не установлено в дипломатическом исходнике.\par
Историческое примечание исходной транскрипции: Оригинал 2408×3448; координаты (x0,y0,x1,y1)=(1830, 2200, 2100, 2350), начало в левом верхнем углу. Просмотр увеличенного фрагмента: 4×. В наборе выбрана ρ; графическое отличие от латинского p в этом почерке недостаточно надёжно.\par\medskip
\par\noindent\leavevmode\phantomsection\label{GELFAND-0154-U009}\textbf{\texttt{GELFAND-0154-U009}} (eq2-variable).\nopagebreak[4]\par
Под второй суммой формулы (2) видна ξ, а в аргументе φ(T\textasciicircum{}k x) — x.\par
Кандидаты: ξ∈Fix T\textasciicircum{}n; φ(T\textasciicircum{}k x). Чтение не установлено в дипломатическом исходнике.\par
Историческое примечание исходной транскрипции: Оригинал 2408×3448; координаты (x0,y0,x1,y1)=(965, 2490, 1100, 2640), начало в левом верхнем углу. Просмотр увеличенного фрагмента: 5×. Несогласованность переменной сохранена, без замены ξ на x или x на ξ.\par\medskip
\subsection*{\texttt{GELFAND-P001417}}
\par\noindent\leavevmode\phantomsection\label{GELFAND-0154-U010}\textbf{\texttt{GELFAND-0154-U010}} (continuation-question-overwriting).\nopagebreak[4]\par
В вопросе о продолжении густо обведены конец сокращения «Сущ.» и начальная буква следующего слова «аналитическое». Остальной вопрос читается.\par
Кандидаты: Сущ. ли аналитическое продолжение?. Чтение не установлено в дипломатическом исходнике.\par
Историческое примечание исходной транскрипции: Оригинал 2408×3448; координаты (x0,y0,x1,y1)=(190, 2005, 680, 2145), начало в левом верхнем углу. Просмотр увеличенного фрагмента: 3×. Поздний слой читается; первоначальные буквы под обводкой не определены. Не подразумевается восстановление стёртого текста.\par\medskip
\subsection*{\texttt{GELFAND-P001418}}
\par\noindent\leavevmode\phantomsection\label{GELFAND-0154-U011}\textbf{\texttt{GELFAND-0154-U011}} (basin-cancelled-sign).\nopagebreak[4]\par
Перед словами «доля воды» расположен одиночный перечёркнутый знак с верхним и нижним штрихами.\par
Кандидаты: не записаны. Чтение не установлено в дипломатическом исходнике.\par
Историческое примечание исходной транскрипции: Оригинал 2408×3448; координаты (x0,y0,x1,y1)=(580, 245, 1830, 365), начало в левом верхнем углу. Просмотр увеличенного фрагмента: 3×. Не установлено, была ли здесь буква или начатый математический знак; отменённый след сохранён отдельным маркером.\par\medskip
\par\noindent\leavevmode\phantomsection\label{GELFAND-0154-U012}\textbf{\texttt{GELFAND-0154-U012}} (eq7-upper-product-bound).\nopagebreak[4]\par
В (7) нижний предел произведения ℓ=0, верхний — n без −1.\par
Кандидаты: ∏\_\{ℓ=0\}\textasciicircum{}\{n\}. Чтение не установлено в дипломатическом исходнике.\par
Историческое примечание исходной транскрипции: Оригинал 2408×3448; координаты (x0,y0,x1,y1)=(1240, 1750, 1450, 2000), начало в левом верхнем углу. Просмотр увеличенного фрагмента: 4×. Возможная ошибка источника: в соответствующей формуле (2) на кадре GELFAND-P001416 предел равен n−1; здесь получается n+1 множителей. Верхний предел n сохранён без исправления.\par\medskip
\par\noindent\leavevmode\phantomsection\label{GELFAND-0154-U013}\textbf{\texttt{GELFAND-0154-U013}} (before-m-prime-cancelled).\nopagebreak[4]\par
Между F′=F\textasciicircum{}k и M′⊂F′×F′ несколько знаков закрыты косой штриховкой; различимы отдельные петли и штрихи, но не полная запись.\par
Кандидаты: не записаны. Чтение не установлено в дипломатическом исходнике.\par
Историческое примечание исходной транскрипции: Оригинал 2408×3448; координаты (x0,y0,x1,y1)=(780, 2895, 970, 3030), начало в левом верхнем углу. Просмотр увеличенного фрагмента: 4×. Содержание отменённой записи не восстанавливалось по следующему определению.\par\medskip
\par\noindent\leavevmode\phantomsection\label{GELFAND-0154-U014}\textbf{\texttt{GELFAND-0154-U014}} (eq10-overwritten-indices).\nopagebreak[4]\par
В (10) малые индексы у второго кортежа переписаны поверх прежних цифр. Поздний слой похож на 2 и k+1; последняя единица в i+1 усилена.\par
Кандидаты: (ξ\_2,…,ξ\_\{k+1\}); (ξ\_i,ξ\_\{i+1\})∈M. Чтение не установлено в дипломатическом исходнике.\par
Историческое примечание исходной транскрипции: Оригинал 2408×3448; координаты (x0,y0,x1,y1)=(990, 3040, 1520, 3190), начало в левом верхнем углу. Просмотр увеличенного фрагмента: 3×. Поздний слой использован в наборе. Первоначальные индексы под обводкой не определены; сдвиговая интерпретация не принимается за доказательство чтения.\par\medskip
\subsection*{\texttt{GELFAND-P001419}}
\par\noindent\leavevmode\phantomsection\label{GELFAND-0154-U015}\textbf{\texttt{GELFAND-0154-U015}} (nb-cancelled-letter).\nopagebreak[4]\par
Между «делить,» и «извлекать» виден короткий густой наклонный след без различимых внутренних букв.\par
Кандидаты: не записаны. Чтение не установлено в дипломатическом исходнике.\par
Историческое примечание исходной транскрипции: Оригинал 2408×3448; координаты (x0,y0,x1,y1)=(1370, 122, 1478, 231), начало в левом верхнем углу. Просмотр увеличенного фрагмента: 4×. Отмечен как отменённый/залитый знак, без догадки о его букве.\par\medskip
\par\noindent\leavevmode\phantomsection\label{GELFAND-0154-U016}\textbf{\texttt{GELFAND-0154-U016}} (heading-last-sign).\nopagebreak[4]\par
После «физики» стоит жирный наклонный штрих, вверху с раздвоением, и отдельная нижняя точка.\par
Кандидаты: !; ?. Чтение не установлено в дипломатическом исходнике.\par
Историческое примечание исходной транскрипции: Оригинал 2408×3448; координаты (x0,y0,x1,y1)=(2238, 424, 2315, 510), начало в левом верхнем углу. Просмотр увеличенного фрагмента: 5×. Вероятнее восклицательный знак; из-за утолщения верхнего штриха вопросительный знак не исключён.\par\medskip
\par\noindent\leavevmode\phantomsection\label{GELFAND-0154-U017}\textbf{\texttt{GELFAND-0154-U017}} (faltings-name).\nopagebreak[4]\par
Фамилия в строке «Виттен слушал доклад …» читается как «Фалтинса»; самостоятельная г после н не выделяется.\par
Кандидаты: Фалтинса; Фалтингса. Чтение не установлено в дипломатическом исходнике.\par
Историческое примечание исходной транскрипции: Оригинал 2408×3448; координаты (x0,y0,x1,y1)=(915, 558, 1420, 668), начало в левом верхнем углу. Просмотр увеличенного фрагмента: 3×. Основной кандидат сохранён буквально. Ни написание, ни отождествление с историческим лицом не исправлялись по внешним сведениям.\par\medskip
\par\noindent\leavevmode\phantomsection\label{GELFAND-0154-U018}\textbf{\texttt{GELFAND-0154-U018}} (witten-list-final-sign).\nopagebreak[4]\par
Справа от «и кто-то ещё» отдельный высокий знак с верхней петлёй и пересекающим нижним росчерком.\par
Кандидаты: \&. Чтение не установлено в дипломатическом исходнике.\par
Историческое примечание исходной транскрипции: Оригинал 2408×3448; координаты (x0,y0,x1,y1)=(1240, 1655, 1360, 1785), начало в левом верхнем углу. Просмотр увеличенного фрагмента: 5×. Нет основания считать знак именем соавтора или восстанавливать имя. Возможная отмена предыдущего знака также не исключена.\par\medskip
\par\noindent\leavevmode\phantomsection\label{GELFAND-0154-U019}\textbf{\texttt{GELFAND-0154-U019}} (eq3-pi-exponent).\nopagebreak[4]\par
Над π в (3) написано −s+1; дробная черта и знаменатель 2 отсутствуют.\par
Кандидаты: π\textasciicircum{}\{−s+1\}. Чтение не установлено в дипломатическом исходнике.\par
Историческое примечание исходной транскрипции: Оригинал 2408×3448; координаты (x0,y0,x1,y1)=(407, 2505, 674, 2635), начало в левом верхнем углу. Просмотр увеличенного фрагмента: 5×. При формальном преобразовании записанного выше (2) получается показатель 1/2−s. Это математическая проверка, а не восстановление источника; в body.tex сохранён −s+1.\par\medskip
\subsection*{\texttt{GELFAND-P001420}}
\par\noindent\leavevmode\phantomsection\label{GELFAND-0154-U020}\textbf{\texttt{GELFAND-0154-U020}} (normalization-cancelled).\nopagebreak[4]\par
После «Нормируем» короткая запись закрыта несколькими косыми штрихами; далее читается «за счёт целочисл. матриц».\par
Кандидаты: не записаны. Чтение не установлено в дипломатическом исходнике.\par
Историческое примечание исходной транскрипции: Оригинал 2408×3448; координаты (x0,y0,x1,y1)=(620, 555, 750, 657), начало в левом верхнем углу. Просмотр увеличенного фрагмента: 4×. Отменённая запись не подменена предполагаемым местоимением.\par\medskip
\par\noindent\leavevmode\phantomsection\label{GELFAND-0154-U021}\textbf{\texttt{GELFAND-0154-U021}} (margin-cancelled).\nopagebreak[4]\par
В правой маргиналии после gl(n,R) находится перечёркнутый знак; далее «станд. базис».\par
Кандидаты: не записаны. Чтение не установлено в дипломатическом исходнике.\par
Историческое примечание исходной транскрипции: Оригинал 2408×3448; координаты (x0,y0,x1,y1)=(2120, 651, 2190, 750), начало в левом верхнем углу. Просмотр увеличенного фрагмента: 5×. Исходная буква или пунктуация не различается.\par\medskip
\par\noindent\leavevmode\phantomsection\label{GELFAND-0154-U022}\textbf{\texttt{GELFAND-0154-U022}} (margin-volume-formula).\nopagebreak[4]\par
В рамке на правом поле записано отношение объёмных форм. Знаменатель читается d tr; малые буквы после vol похожи на sl и gl.\par
Кандидаты: vol\_sl = vol\_gl / d tr. Чтение не установлено в дипломатическом исходнике.\par
Историческое примечание исходной транскрипции: Оригинал 2408×3448; координаты (x0,y0,x1,y1)=(1980, 710, 2320, 853), начало в левом верхнем углу. Просмотр увеличенного фрагмента: 4×. Сравнение с предыдущими словами об алгебрах Ли поддерживает чтение sl/gl. Знаменатель не заменялся на det или другую стандартную формулу.\par\medskip
\par\noindent\leavevmode\phantomsection\label{GELFAND-0154-U023}\textbf{\texttt{GELFAND-0154-U023}} (invariant-volume-cancelled).\nopagebreak[4]\par
Между «есть» и vol\_\{H²\}⊗dφ находится густо зачернённый короткий вертикально-наклонный знак.\par
Кандидаты: не записаны. Чтение не установлено в дипломатическом исходнике.\par
Историческое примечание исходной транскрипции: Оригинал 2408×3448; координаты (x0,y0,x1,y1)=(1600, 990, 1690, 1120), начало в левом верхнем углу. Просмотр увеличенного фрагмента: 4×. Перечёркнутый след учтён отдельно; не принят за интеграл или множитель.\par\medskip
\subsection*{\texttt{GELFAND-P001421}}
\par\noindent\leavevmode\phantomsection\label{GELFAND-0154-U024}\textbf{\texttt{GELFAND-0154-U024}} (witten-heading-faint-word).\nopagebreak[4]\par
Между «Пример» и «работ Виттена» почти выцветшее слово; отдельные следы похожи на «силы».\par
Кандидаты: силы. Чтение не установлено в дипломатическом исходнике.\par
Историческое примечание исходной транскрипции: Оригинал 2408×3448; координаты (x0,y0,x1,y1)=(600, 1395, 837, 1520), начало в левом верхнем углу. Просмотр увеличенного фрагмента: 4×. Чтение вероятное, но потеря штрихов значительна; слово не дополнялось по внешней формулировке.\par\medskip
\par\noindent\leavevmode\phantomsection\label{GELFAND-0154-U025}\textbf{\texttt{GELFAND-0154-U025}} (book-title-cancelled).\nopagebreak[4]\par
В названии книги после «Спектры,» старое слово густо перечёркнуто, а над ним со знаком вставки написано «сигналы».\par
Кандидаты: не записаны. Чтение не установлено в дипломатическом исходнике.\par
Историческое примечание исходной транскрипции: Оригинал 2408×3448; координаты (x0,y0,x1,y1)=(1650, 1938, 1975, 2060), начало в левом верхнем углу. Просмотр увеличенного фрагмента: 3×. Поздняя вставка читается; прежнее слово не восстановлено.\par\medskip
\par\noindent\leavevmode\phantomsection\label{GELFAND-0154-U026}\textbf{\texttt{GELFAND-0154-U026}} (surface-symbol).\nopagebreak[4]\par
В строке «Рассм. поверхность» стоит каллиграфическая F, та же форма буквы затем встречается во втором аргументе H\textasciicircum{}i(X,F).\par
Кандидаты: поверхность F; H\textasciicircum{}i(X,F). Чтение не установлено в дипломатическом исходнике.\par
Историческое примечание исходной транскрипции: Оригинал 2408×3448; координаты (x0,y0,x1,y1)=(802, 2240, 1435, 2355), начало в левом верхнем углу. Просмотр увеличенного фрагмента: 3×. Не заменено на X. Возможное смешение обозначения поверхности и коэффициентов/пучка остаётся свойством источника; новая интерпретация в транскрипцию не введена.\par\medskip
\par\noindent\leavevmode\phantomsection\label{GELFAND-0154-U027}\textbf{\texttt{GELFAND-0154-U027}} (before-vector-spaces-cancelled).\nopagebreak[4]\par
В скобках перед «вект. пр-ва» зачёркнуто слово с несколькими видимыми верхними штрихами и окончанием, похожим на «ные».\par
Кандидаты: конечные. Чтение не установлено в дипломатическом исходнике.\par
Историческое примечание исходной транскрипции: Оригинал 2408×3448; координаты (x0,y0,x1,y1)=(488, 2720, 795, 2840), начало в левом верхнем углу. Просмотр увеличенного фрагмента: 4×. Слово «конечные» — кандидат, не уверенная расшифровка.\par\medskip
\subsection*{\texttt{GELFAND-P001422}}
\par\noindent\leavevmode\phantomsection\label{GELFAND-0154-U028}\textbf{\texttt{GELFAND-0154-U028}} (cancelled-divisor-map).\nopagebreak[4]\par
Между (3) и заголовком «Ф-ла Делиня» зачёркнута короткая формула: два знака, сходных с D, стрелка и конечное K; промежуточные штрихи пересекаются с зачёркиванием.\par
Кандидаты: D → D+K; D ↦ D+K; D → D⊗K. Чтение не установлено в дипломатическом исходнике.\par
Историческое примечание исходной транскрипции: Оригинал 2408×3448; координаты (x0,y0,x1,y1)=(325, 880, 740, 1000), начало в левом верхнем углу. Просмотр увеличенного фрагмента: 4×. Кандидаты передают лишь различимые очертания; операцию и вид стрелки уверенно восстановить нельзя. Формула не подменена ожидаемой математической операцией.\par\medskip
\par\noindent\leavevmode\phantomsection\label{GELFAND-0154-U029}\textbf{\texttt{GELFAND-0154-U029}} (eq4-cancelled-sign).\nopagebreak[4]\par
Справа от равенства (4) виден одиночный высокий знак, перечёркнутый несколькими короткими косыми штрихами.\par
Кандидаты: t; неразборчивый зачёркнутый знак. Чтение не установлено в дипломатическом исходнике.\par
Историческое примечание исходной транскрипции: Оригинал 2408×3448; координаты (x0,y0,x1,y1)=(957, 1110, 1080, 1260), начало в левом верхнем углу. Просмотр увеличенного фрагмента: 4×. Знак не присоединён к формуле как определённый символ.\par\medskip
\par\noindent\leavevmode\phantomsection\label{GELFAND-0154-U030}\textbf{\texttt{GELFAND-0154-U030}} (eq4-determinant-convention).\nopagebreak[4]\par
В (4) явно написано det H\textasciicircum{}• = det F\_K; отрицательный показатель справа отсутствует. Выше записано 2g−2=deg K; на предыдущем кадре дана конвенция det H\textasciicircum{}•=det H\textasciicircum{}0 det H\textasciicircum{}2/det H\textasciicircum{}1.\par
Кандидаты: det H\textasciicircum{}• = det F\_K (видимая запись). Чтение не установлено в дипломатическом исходнике.\par
Историческое примечание исходной транскрипции: Оригинал 2408×3448; координаты (x0,y0,x1,y1)=(265, 1100, 815, 1245), начало в левом верхнем углу. Просмотр увеличенного фрагмента: 3×. Математическая проверка, не исправление: для постоянной системы ранга r записанная на предыдущем кадре конвенция даёт степень 2−2g у определителя слоя, тогда как формула (3) для K имеет степень 2g−2. Возможно, недостаёт обратной степени либо подразумевается иная конвенция. По изображению это не разрешено; исходное равенство сохранено.\par\medskip
\subsection*{\texttt{GELFAND-P001423}}
\par\noindent\leavevmode\phantomsection\label{GELFAND-0154-U031}\textbf{\texttt{GELFAND-0154-U031}} (theorem-cancelled-sign).\nopagebreak[4]\par
В первой строке теоремы между «пр-в» и E расположен плотный зачёркнутый знак.\par
Кандидаты: с; неразборчивый зачёркнутый знак. Чтение не установлено в дипломатическом исходнике.\par
Историческое примечание исходной транскрипции: Оригинал 2408×3448; координаты (x0,y0,x1,y1)=(2090, 535, 2260, 680), начало в левом верхнем углу. Просмотр увеличенного фрагмента: 5×. Буква не восстановлена по грамматически ожидаемому обороту. Слово «набором» на следующей строке сохранено без добавления неустановленного предлога.\par\medskip
\par\noindent\leavevmode\phantomsection\label{GELFAND-0154-U032}\textbf{\texttt{GELFAND-0154-U032}} (representation-cancelled-sign).\nopagebreak[4]\par
В записи α∈Σ перед Σ зачёркнут небольшой знак; при увеличении видна перекрёстная штриховка.\par
Кандидаты: вертикальная черта; неразборчивый зачёркнутый знак. Чтение не установлено в дипломатическом исходнике.\par
Историческое примечание исходной транскрипции: Оригинал 2408×3448; координаты (x0,y0,x1,y1)=(1995, 1570, 2210, 1780), начало в левом верхнем углу. Просмотр увеличенного фрагмента: 5×. Возможно начало обозначения |Σ|, однако математическая аналогия с теоремой выше не устанавливает чтение.\par\medskip
\par\noindent\leavevmode\phantomsection\label{GELFAND-0154-U033}\textbf{\texttt{GELFAND-0154-U033}} (tetrahedron-cancelled-space).\nopagebreak[4]\par
После «таков: в» перед E, dim E=3 зачёркнуто обозначение пространства и сокращение; различим начальный знак, похожий на R, верхний показатель и окончание «пр-ве».\par
Кандидаты: R\textasciicircum{}3-пр-ве; R\textasciicircum{}3 пр-ве; R\textasciicircum{}3. Чтение не установлено в дипломатическом исходнике.\par
Историческое примечание исходной транскрипции: Оригинал 2408×3448; координаты (x0,y0,x1,y1)=(615, 2975, 1045, 3145), начало в левом верхнем углу. Просмотр увеличенного фрагмента: 5×. Сочетание частично перекрыто линиями зачёркивания. Незачёркнутая последующая запись E, dim E=3 воспроизведена независимо от кандидата.\par\medskip
\subsection*{\texttt{GELFAND-P001424}}
\par\noindent\leavevmode\phantomsection\label{GELFAND-0154-U034}\textbf{\texttt{GELFAND-0154-U034}} (kac-initial).\nopagebreak[4]\par
В начале строки «У В. … Каца» второй инициал густо перезаписан и зачёркнут; отдельная В перед ним читается.\par
Кандидаты: Г.; Д.; неразборчивый исправленный инициал. Чтение не установлено в дипломатическом исходнике.\par
Историческое примечание исходной транскрипции: Оригинал 2408×3448; координаты (x0,y0,x1,y1)=(155, 770, 435, 920), начало в левом верхнем углу. Просмотр увеличенного фрагмента: 7×. Имя не дополнено по сведениям об известном математике. Переданы читаемая В и факт исправления второго инициала.\par\medskip
\par\noindent\leavevmode\phantomsection\label{GELFAND-0154-U035}\textbf{\texttt{GELFAND-0154-U035}} (cancelled-seminar-line).\nopagebreak[4]\par
Перед «След. раз: 8 февраля» целиком зачёркнута строка. Начало похоже на «Семинар», затем несколько слов под длинными горизонтальными и косыми штрихами.\par
Кандидаты: Семинар закрывается …; Семинар …. Чтение не установлено в дипломатическом исходнике.\par
Историческое примечание исходной транскрипции: Оригинал 2408×3448; координаты (x0,y0,x1,y1)=(235, 1300, 1250, 1430), начало в левом верхнем углу. Просмотр увеличенного фрагмента: 5×. Начальное слово также передано в составе неопределённого чтения. Окончание не восстановлено из даты следующего заседания; смысл зачёркнутого объявления не установлен.\par\medskip
\par\noindent\leavevmode\phantomsection\label{GELFAND-0154-U036}\textbf{\texttt{GELFAND-0154-U036}} (cancelled-seminar-content).\nopagebreak[4]\par
Густо зачёркнутая строка между «(докл. на конференции)» и «След. раз: 8 февраля» не поддаётся полному чтению; в основном тексте оставлено только неопределённое начальное слово.\par
Кандидаты: не записаны. Чтение не установлено в дипломатическом исходнике.\par
Историческое примечание исходной транскрипции: Оригинал 2408×3448; координаты (x0,y0,x1,y1)=(200, 1280, 1310, 1410), начало в левом верхнем углу. Просмотр увеличенного фрагмента: 7×. Весь кадр доступен и осмотрен. Неразборчивость касается именно этой зачёркнутой строки, а не остального текста. Повторный просмотр фрагмента с увеличением до 7× не установил полного чтения. Кандидаты отдельных слов вынесены в GELFAND-0154-U035. Эта запись явно учитывает фактически невосстановленный текст.\par\medskip
